\section{Discusi\'on}
\label{sec:discusion}

\subsection{Por qu\'e el piloto enga\~na}

La ca\'ida de Recall@10 de 0,947 a 0,684 al pasar del piloto al corpus completo (Tabla~\ref{tab:e6}) no es un fallo de configuraci\'on: es el efecto de la competencia. Con 1.932 documentos, el fragmento relevante compite contra pocos vecinos; con 161.868, compite contra cientos de fragmentos que comparten vocabulario gen\'erico. El texto sin cierre terminal (32,7\,\%) agrava el problema cuando el corte elimina la frase que responde, y el chunking fijo puede partir la evidencia en dos fragmentos que se disputan el mismo puesto. Ninguna de esas causas se puede medir con una muestra peque\~na. La consecuencia pr\'actica es que un piloto sirve para depurar el pipeline, no para elegir codificador, chunking o $k$. Esa elecci\'on hay que hacerla a la escala final, aunque cueste horas de GPU.

\subsection{El coste del orden}

El reordenador es una herramienta de precisi\'on, no de cobertura. Sobre el top-5 sube el nDCG 5,6 puntos y el P@1 13,1 con 1,30 s de coste, porque el recuperador ya trajo la respuesta y solo hab\'ia que colocarla primero. Sobre el top-20 baja el recall y el P@5: al ampliar el conjunto entran fragmentos peores que desplazan a los buenos. El reordenador no puede recuperar lo que el recuperador no trajo, y por eso no arregla g-046 ni g-031. La decisi\'on de adoptarlo depende de d\'onde duela el error: si el generador se pierde con la evidencia en el medio del contexto \cite{liu2024lost}, ordenar importa; si el problema es que falta el documento, ordenar no sirve.

\subsection{La abstenci\'on barata}

La compuerta cuesta milisegundos y rechaza las 5 preguntas sin respuesta; el juez cuesta una generaci\'on completa. Esa asimetr\'ia justifica el clasificador. El precio es un falso rechazo en 0,494, una pregunta leg\'itima sobre la inflaci\'on que queda fuera por seis mil\'esimas. El umbral es una decisi\'on de negocio, no t\'ecnica: aceptar una pregunta ajena erosiona la confianza m\'as que abstenerse en una propia. En un asistente de corpus, el usuario prefiere un "no lo s\'e" a una respuesta inventada, y por eso se mantiene 0,5 aunque el test mejore con 0,45. Si el coste de abstenerse fuera mayor, el umbral se mover\'ia en la otra direcci\'on.

\subsection{La evaluaci\'on}

La precisi\'on del contexto de RAGAS pas\'o de 0,750 a 0,881 en dos corridas del mismo sistema. La primera cifra ten\'ia 4 valores de 38 porque el juez no terminaba; la segunda tiene 34. La conclusi\'on anterior, que atribu\'ia a la arquitectura una p\'erdida de precisi\'on, era un artefacto del evaluador. La comparaci\'on con la arquitectura previa (0,760, con 36 valores) es otra cosa: cambia juez y embedding, as\'i que se reporta como observaci\'on y no como mejora causal. El episodio deja dos reglas. La primera: reportar $n$ junto a cada media y no comparar m\'etricas con cobertura distinta. La segunda: el ruido del juez no es aleatorio, se concentra en contextos largos y en fallos de parser, as\'i que un timeout m\'as generoso y menos concurrencia lo reducen.

El acuerdo humano tambi\'en delimita lo que se puede afirmar. La correcci\'on (0,95) y las citas (0,90) son fiables; la abstenci\'on (0,72) es aceptable; la fidelidad (0,12) no. Usar la fidelidad para elegir modelo habr\'ia sido elegir con ruido. La r\'ubrica de fidelidad mezcla errores factuales con abstenciones injustificadas, y los anotadores no comparten el criterio. Un eje que no alcanza acuerdo no debe decidir, y eso es una lecci\'on de dise\~no de la evaluaci\'on, no un defecto de los anotadores.

\subsection{Lo que mostr\'o el conjunto reservado}

El conjunto reservado cambia el diagn\'ostico. En el gold exploratorio el recall@10 era 0,912 y el sistema entregaba 37 de 38; en 20 preguntas nuevas el recall baja a 0,600 y entrega 17. La ca\'ida no se debe al azar: ocho preguntas no recuperan el documento de referencia y tres de ellas terminan en abstenci\'on. El sistema sabe decir que no, pero tambi\'en responde con una cifra distinta cuando el pasaje correcto no est\'a (r-006). Ese es el error que m\'as importa en un asistente de prensa: una respuesta con cita y con el n\'umero equivocado.

La compuerta anterior tampoco generalizaba el cero: rechazaba 4 de 5 negativas nuevas y dejaba entrar una con 0,579, que el guardarra\'il deten\'ia despu\'es. La compuerta nueva, con la puntuaci\'on del cross-encoder, rechaza las 5 y baja los falsos rechazos de 6 a 3. Los tres que quedan son dos fallos de recuperaci\'on (r-012 y r-014) y un falso rechazo con evidencia recuperada (r-011). El 0/9 del conjunto anterior era una muestra demasiado peque\~na para sostener una tasa nula. La lecci\'on no es retirar la compuerta, sino reportar su tasa con el denominador y no prometer ausencia de falsos aceptados.

El conjunto reservado tambi\'en muestra que el gold exploratorio no era representativo en idioma: en ingl\'es el recall es 0,625 y el nDCG 0,579; en espa\~nol, 0,583 y 0,409. El ingl\'es, que en el gold exploratorio parec\'ia el punto d\'ebil, recupera mejor en la muestra nueva. Con 8 y 12 preguntas, las diferencias no son concluyentes, pero bastan para no tratar el idioma como una conclusi\'on cerrada.

\subsection{Qu\'e se llevar\'ia a producci\'on}

La configuraci\'on que se llevar\'ia es la de la Tabla~\ref{tab:config-final}: bge-m3 con \texttt{ef}=800, denso, $k=5$, reranker en CPU, gemma4:e4b con MTP, prompt v2 y compuerta con cross-encoder y umbral 0,4. Es la que mejor equilibra cobertura, orden, latencia y coste en este corpus. Las dos deudas t\'ecnicas son la recuperaci\'on de g-031 y g-046 y la correcci\'on de la respuesta. La primera exige un conjunto de candidatos mayor o una consulta distinta; la segunda, un generador mejor o un prompt que aproveche m\'as la evidencia. Ninguna se resuelve con m\'as reordenamiento.

El monitoreo que acompa\~nar\'ia el despliegue tiene cuatro se\~nales: recall en una muestra fija del gold, tasa de abstenci\'on, tasa de citas v\'alidas y latencia p95. Las cuatro se pueden calcular con la infraestructura actual y detectan las regresiones que importan: que el \'indice cambie, que el modelo se actualice, que el prompt se rompa o que el hardware se quede sin memoria. La compuerta se monitorea aparte, con un conjunto fijo de preguntas de fuera, porque su deriva no se ve en las m\'etricas de respuesta.
